\documentclass[12pt]{article}

\usepackage{amsmath}
\usepackage{amssymb}
\usepackage{physics}
\usepackage{hyperref}

\title{Affine Geometry}
\author{}
\date{}

\begin{document}

\maketitle

% ============================================================
\section{Introduction}
% ============================================================

\section{Affine Spaces}

\textbf{Definition of an affine space.}

An \(n\)-dimensional affine space over the field of real numbers
\(\mathbb{R}\) consists of:

\begin{itemize}
    \item a vector space \(V\), whose elements are called vectors;
    \item a set \(X\), whose elements are called points.
\end{itemize}

To each vector \(\mathbf{v} \in V\) is associated a mapping
\[
\mathbf{v} \colon X \longrightarrow X,
\]
\[
\mathbf{v}(x) = y,
\]
where \(x,y \in X\), with the following properties:

\begin{enumerate}
    \item
    \(\mathbf{0}(x)=x\), where \(\mathbf{0}\) is the null vector.

    \item
    For all \(x,y\in X\), there exists a unique \(\mathbf{v}\in V\)
    such that \(\mathbf{v}(x)=y\), and we denote
    \(\mathbf{v}=\overrightarrow{xy}\).

    \item
    If \(\mathbf{v},\mathbf{w}\in V\), the composition of the action
    of two vectors is defined formally as the sum of two vectors:
    
    \[(\mathbf{v}\circ\mathbf{w})(x)
    =\mathbf{v}\bigl(\mathbf{w}(x)\bigr)
    \coloneqq(\mathbf{v}+\mathbf{w})(x).\]

\end{enumerate}

Such a composition must be commutative, as the sum of vectors is:
\[
\mathbf{v}\circ\mathbf{w}
=
\mathbf{v}+\mathbf{w}
=
\mathbf{w}+\mathbf{v}
=
\mathbf{w}\circ\mathbf{v}.
\] 

\begin{theorem}\label{thm:affine-composition}
The composition of translations is commutative.
\end{theorem}

By Theorem \ref{thm:affine-composition}, the composition is commutative.

\end{document}
